% Figura 3 del informe corto: las piezas que existen y la que falta.
% Fuente: docs/investigacion/10_problema_usuarios_alternativas.md, secciones 1 a 3.
% Se incluye con \input desde informe_corto.tex, que define los colores tinta, ambar, ceniza, hielo y papel.
\begin{tikzpicture}[
  font=\small,
  pieza/.style={draw=tinta, fill=hielo, line width=0.6pt, text width=3.2cm, minimum height=1.75cm,
    align=center, inner sep=4pt},
  falta/.style={draw=ambar, fill=papel, dashed, line width=0.9pt, text width=10.2cm, align=center,
    inner sep=6pt},
]
\node[pieza] (rud) at (0,0)
  {{\color{tinta}\textsc{rud}}\\{\footnotesize UNGRD}\\[3pt]{\footnotesize\color{ceniza} quién es damnificado}};
\node[pieza] (kobo) at (4.0,0)
  {{\color{tinta}KoboCollect y AccessRC}\\{\footnotesize Cruz Roja}\\[3pt]{\footnotesize\color{ceniza} registro de hogares}};
\node[pieza] (suma) at (8.0,0)
  {{\color{tinta}\textsc{suma}/\textsc{lss}}\\{\footnotesize OPS}\\[3pt]{\footnotesize\color{ceniza} inventario y entregas}};
\node[pieza] (abaco) at (12.0,0)
  {{\color{tinta}Plataforma de ABACO}\\{\footnotesize red de bancos}\\[3pt]{\footnotesize\color{ceniza} recaudo y reporte al donante}};
\node[falta] (despacho) at (6.0,-2.75)
  {{\color{ambar}\textsc{la pieza que falta: el despacho}}\\[3pt]
   pedido \enspace$\rightarrow$\enspace vehículo \enspace$\rightarrow$\enspace vía \enspace$\rightarrow$\enspace entrega\\[3pt]
   {\footnotesize\color{ceniza}Es lo que propone DespachoX. Una persona aprueba cada plan.}};
\draw[-{Stealth[length=2.6mm]}, ambar, line width=0.9pt]
  (despacho.north -| suma.south) -- node[right, font=\footnotesize, text=ambar] {exporta} (suma.south);
\end{tikzpicture}
